% ============================ 第二节 概览 ====================================
% 结构对照准则第 57 号第二十至二十二条：显要位置声明（第二十一条原句）→
% 一、重大事项提示（第二十二条第（一）项：重大风险扼要披露并索引至相关章节）
% → 二、发行人及中介机构基本情况（合一四列表＋通栏小标题）→ 三、本次发行概况
% （参考格式：基本情况+重要日期）→ 主营业务 → 科创板定位 → 财务数据 →
% 上市标准 → 治理特殊安排 → 募集资金用途。数据与第五、六、七节完全勾稽。
\gssection{概览}

\zd{本概览仅对招股说明书全文作扼要提示。投资者作出投资决策前，应认真阅读本
招股说明书全文。}

\subsection{重大事项提示}

\subsubsection{特别风险提示}

本公司特别提请投资者认真阅读本招股说明书``第三节\hspace{0.5\ccwd}风险因素''
的全部内容，并特别关注下列风险：

\paragraph{算力供应及核心芯片受限风险}
大模型的训练与推理依赖高性能智能计算芯片及配套算力基础设施。受国际贸易环境
及相关出口管制政策影响，高性能芯片的供应存在不确定性；若算力供应紧张或采购
成本大幅上升，将对公司模型迭代进度、服务质量及盈利能力产生不利影响。详见
``第三节\hspace{0.5\ccwd}风险因素''之``二、（一）''。

\paragraph{技术升级与产品迭代风险}
人工智能大模型技术仍处于快速演进阶段，模型架构、训练方法与产品形态迭代频繁。
若公司对技术路线的判断出现偏差，或者新一代模型的研发进度与效果不及预期，公
司产品和服务的市场竞争力将受到不利影响。详见``第三节\hspace{0.5\ccwd}风险因
素''之``一、（一）''。

\paragraph{客户相对集中风险}
报告期各期，公司对前五大客户的销售收入占营业收入的比例分别为42.56\%、
43.28\%和43.92\%，客户集中度相对较高。若主要客户经营情况发生不利变化或减少
对公司产品及服务的采购，公司经营业绩存在下滑的风险。详见``第三节
\hspace{0.5\ccwd}风险因素''之``一、（四）''。

\paragraph{应收账款回收风险}
报告期各期末，公司应收账款账面余额分别为18,236.45万元、23,845.67万元和
30,124.56万元，占当期营业收入的比例分别为22.12\%、22.66\%和22.72\%。若下游
客户资信状况恶化，公司存在应收账款不能按期收回或无法收回的风险。详见``第三
节\hspace{0.5\ccwd}风险因素''之``一、（五）''。

\subsubsection{本次发行前滚存利润的分配安排}

经公司2026年第一次临时股东会审议通过，本次发行完成前公司滚存的未分配利润，
由本次发行完成后的新老股东按发行后的持股比例共同享有。

\subsubsection{本次发行上市后的股利分配政策}

本次发行上市后，公司将实行持续、稳定的利润分配政策，优先采用现金分红方式。
具体政策详见``第九节\hspace{0.5\ccwd}投资者保护''之``一、本次发行后的股利分
配政策''。

\subsubsection{财务报告审计截止日后的主要财务信息及经营状况}

本招股说明书财务报告审计截止日为2025年12月31日。示例会计师对公司2026年
1—3月的财务报表进行了审阅，并出具了《审阅报告》（示例会所阅字〔2026〕第
××号）。审计截止日后，公司经营模式、主要客户及供应商构成、税收政策等未发
生重大变化，主要经审阅财务数据如下：

\biaodw{万元}
\begin{gstab}
\begin{longtable}{|L{52mm}|R{46mm}|R{40mm}|}
\hline
\thA{项目} & \thBw{46mm}{2026年1—3月／\newline 2026年3月31日} & \thB{变动幅度} \\ \hline
\endfirsthead
\hline
\thA{项目} & \thBw{46mm}{2026年1—3月／\newline 2026年3月31日} & \thB{变动幅度} \\ \hline
\endhead
营业收入 & 35,246.18 & 24.35\% \\ \hline
归属于母公司股东的净利润 & 5,682.34 & 22.18\% \\ \hline
扣除非经常性损益后归属于母公司股东的净利润 & 5,486.21 & 21.57\% \\ \hline
总资产 & 178,456.32 & 5.98\% \\ \hline
归属于母公司股东的权益 & 102,737.61 & 5.85\% \\ \hline
\end{longtable}
\end{gstab}
\tabnote{营业收入及净利润指标的变动幅度为较上年同期变动，总资产及归属于母公
司股东的权益的变动幅度为较2025年12月31日变动；上述2026年1—3月财务数据经审
阅、未经审计。}

\subsubsection{本公司及相关责任主体作出的重要承诺}

公司及其控股股东、实际控制人、董事、高级管理人员、核心技术人员以及本次发行
的保荐人等责任主体，就股份限售安排、持股意向及减持、稳定股价、欺诈发行购回
股份、填补被摊薄即期回报、依法承担赔偿责任等事项作出了承诺，具体内容详见
``第九节\hspace{0.5\ccwd}投资者保护''。\zd{本公司提请投资者认真阅读上述承诺
的具体内容、履行情况及未能履行承诺时的约束措施。}

\subsection{发行人及本次发行的中介机构基本情况}

\begin{gstab}
\begin{longtable}{|L{26mm}|L{40mm}|L{28mm}|L{38mm}|}
\hline
\multicolumn{4}{|c|}{\bfseries （一）发行人基本情况} \\ \hline
\endfirsthead
\hline
\endhead
发行人名称 & \gsCompany & 成立日期 & 2021年3月18日 \\ \hline
注册资本 & 36,000.00万元 & 法定代表人 & 林×× \\ \hline
注册地址 & \gsRegAddr & 主要生产经营地址 & \gsRegAddr \\ \hline
控股股东 & 示例控股有限公司 & 实际控制人 & 林×× \\ \hline
行业分类 & 软件和信息技术服务业（I65） & 在其他交易场所（申请）挂牌或上市的情况 & 无 \\ \hline
\multicolumn{4}{|c|}{\bfseries （二）本次发行的有关中介机构} \\ \hline
保荐人 & \gsSponsor & 主承销商 & \gsSponsor \\ \hline
发行人律师 & 示例市示例律师事务所 & 其他承销机构 & 无 \\ \hline
审计机构 & 示例会计师事务所（特殊普通合伙） & 评估机构 & 不适用 \\ \hline
\multicolumn{2}{|L{\dimexpr 66mm+2\tabcolsep+\arrayrulewidth\relax}|}{发行人与本次发行有关的保荐人、承销机构、证券服务机构及其负责人、高级管理人员、经办人员之间存在的直接或间接的股权关系或其他利益关系} & \multicolumn{2}{L{\dimexpr 66mm+2\tabcolsep+\arrayrulewidth\relax}|}{除保荐人依法设立的相关子公司拟参与本次发行战略配售外，不存在直接或间接的股权关系或其他利益关系} \\ \hline
\multicolumn{4}{|c|}{\bfseries （三）本次发行其他有关机构} \\ \hline
股票登记机构 & 中国证券登记结算有限责任公司上海分公司 & 收款银行 & \gskong{} \\ \hline
\multicolumn{2}{|L{\dimexpr 66mm+2\tabcolsep+\arrayrulewidth\relax}|}{其他与本次发行有关的机构} & \multicolumn{2}{L{\dimexpr 66mm+2\tabcolsep+\arrayrulewidth\relax}|}{无} \\ \hline
\end{longtable}
\end{gstab}

\subsection{本次发行概况}

\subsubsection{本次发行的基本情况}

\begin{gstab}
\begin{longtable}{|L{45mm}|L{100mm}|}
\hline
股票种类／每股面值 & 人民币普通股（A股）／人民币1.00元 \\ \hline
发行股数 & 12,000.00万股，占发行后总股本的25.00\%；全部为发行新股，无股东公开发售股份 \\ \hline
发行后总股本 & 48,000.00万股 \\ \hline
每股发行价格 & 人民币\gskong 元 \\ \hline
发行市盈率 & \gskong 倍（按发行价格除以发行后每股收益计算） \\ \hline
发行前每股净资产 & 2.70元（按2025年12月31日经审计归属于母公司股东的权益除以发行前总股本计算） \\ \hline
发行后每股净资产 & \gskong 元 \\ \hline
发行前每股收益 & 0.61元（按2025年度归属于母公司股东的净利润除以发行前总股本计算） \\ \hline
发行后每股收益 & 0.46元（按2025年度归属于母公司股东的净利润除以发行后总股本计算） \\ \hline
发行市净率 & \gskong 倍（按发行价格除以发行后每股净资产计算） \\ \hline
发行方式 & 网下向符合条件的投资者询价配售与网上向持有上海市场非限售A股股份和非限售存托凭证市值的社会公众投资者定价发行相结合，或中国证监会、上交所认可的其他方式 \\ \hline
发行对象 & 符合资格的询价对象和在上交所开户的境内自然人、法人等投资者（国家法律法规禁止购买者除外） \\ \hline
承销方式 & 余额包销 \\ \hline
募集资金总额 & \gskong 万元 \\ \hline
募集资金净额 & \gskong 万元 \\ \hline
募集资金投资项目 & 新一代大模型研发及训练平台建设项目、智能体应用与行业解决方案研发项目、补充流动资金 \\ \hline
发行费用概算 & \gskong 万元 \\ \hline
保荐人相关子公司拟参与战略配售情况 & 保荐人依法设立的相关子公司拟按照相关规定参与本次发行战略配售 \\ \hline
高级管理人员、员工拟参与战略配售情况 & 无 \\ \hline
\end{longtable}
\end{gstab}

\subsubsection{本次发行上市的重要日期}

\begin{gstab}
\begin{longtable}{|L{45mm}|L{100mm}|}
\hline
刊登发行公告日期 & \gskong 年\gskong 月\gskong 日 \\ \hline
开始询价推介日期 & \gskong 年\gskong 月\gskong 日 \\ \hline
刊登定价公告日期 & \gskong 年\gskong 月\gskong 日 \\ \hline
申购日期和缴款日期 & \gskong 年\gskong 月\gskong 日 \\ \hline
股票上市日期 & \gskong 年\gskong 月\gskong 日 \\ \hline
\end{longtable}
\end{gstab}

\subsection{发行人主营业务经营情况}

公司是一家专注于人工智能大模型的高新技术企业，主营业务为大模型的研发与产业
化应用，向客户提供大模型解决方案、开放平台（MaaS）服务及智能应用产品。公司
自主研发了通用大模型及多模态模型体系，面向金融、政务、制造、互联网等行业客
户提供私有化部署与联合建模服务，并通过开放平台以API方式对外输出模型能力。
报告期各期，公司营业收入分别为82,456.32万元、105,238.47万元和132,594.18万
元，年均复合增长率为26.81\%；归属于母公司股东的净利润分别为12,035.28万元、
16,842.51万元和21,976.35万元，经营业绩持续快速增长。2025年度，大模型解决方
案、开放平台（MaaS）服务、智能应用及其他占主营业务收入的比例分别为55.21\%、
31.27\%和13.52\%，销售以直销模式为主。

\subsection{发行人符合科创板定位}

\subsubsection{公司符合行业领域要求}

根据《科创属性评价指引（试行）（2024年修订）》及《上海证券交易所科创板企业
发行上市申报及推荐暂行规定（2024年4月修订）》等有关规定，公司所属行业领域
情况如下：

\begin{gstab}
% 整表为一个不可分割块（普通 tabular，思朗式真网格）：中列七个勾选项各占
% 一行、\cline{2-2} 分隔线贯穿中列；标签与说明列用 \multirow 跨七行合并；
% 末行撑高（\rule 深度）使总高容纳说明文字——说明文字增删后微调该撑高值
% 与两个 \multirow 的 [垂直偏移]。编译日志中本表产生一条 Overfull \vbox
% 提示，系 multirow 对跨行盒高的估算机制所致，版面输出无损（已逐像素核验）。
{\centering
\begin{tabular}{|>{\centering\arraybackslash}p{26mm}|L{40mm}|>{\raggedright\arraybackslash}p{64mm}|}
\hline
\multirow{7}{=}[-37pt]{\centering 公司所属行业领域} & \zd{\kxv 新一代信息技术} & \multirow{7}{=}[0pt]{公司主营业务为人工智能大模型的研发与产业化应用。根据中国证监会《上市公司行业分类指引》，公司属于``软件和信息技术服务业（I65）''；根据国家统计局《国民经济行业分类》（GB/T 4754—2017），公司属于``I65 软件和信息技术服务业''；根据《战略性新兴产业分类（2018）》，公司属于``1 新一代信息技术产业''中的人工智能相关领域；根据《上海证券交易所科创板企业发行上市申报及推荐暂行规定（2024年4月修订）》，公司属于``新一代信息技术领域''中的``人工智能''领域。} \\ \cline{2-2}
 & \kx 高端装备 & \\ \cline{2-2}
 & \kx 新材料 & \\ \cline{2-2}
 & \kx 新能源 & \\ \cline{2-2}
 & \kx 节能环保 & \\ \cline{2-2}
 & \kx 生物医药 & \\ \cline{2-2}
 & \kx 符合科创板定位的\newline\hspace*{1.05em}其他领域\rule[-120pt]{0pt}{120pt} & \\ \hline
\end{tabular}\par}
\end{gstab}

\subsubsection{公司符合科创属性相关指标要求}

公司符合《科创属性评价指引（试行）（2024年修订）》规定的科创属性评价标准一
的全部指标，具体情况如下：

\begin{gstab}
\begin{longtable}{|L{50mm}|C{20mm}|L{61mm}|}
\hline
\thA{科创属性评价标准一} & \thB{是否符合} & \thB{指标情况} \\ \hline
\endfirsthead
\hline
\thA{科创属性评价标准一} & \thB{是否符合} & \thB{指标情况} \\ \hline
\endhead
最近3年累计研发投入占最近3年累计营业收入比例不低于5\%，或最近3年累计研发投入金额不低于8,000万元 & \mbox{\kxv 是\hspace{0.4\ccwd}\kx 否} & 2023—2025年度，公司累计研发投入为30,128.72万元，占同期累计营业收入的比例为9.41\%，研发投入占比和金额均满足条件 \\ \hline
研发人员占当年员工总数的比例不低于10\% & \mbox{\kxv 是\hspace{0.4\ccwd}\kx 否} & 截至2025年12月31日，公司研发与技术人员726人，占员工总数的比例为56.45\%，满足条件 \\ \hline
应用于公司主营业务并能够产业化的发明专利7项以上 & \mbox{\kxv 是\hspace{0.4\ccwd}\kx 否} & 截至2025年12月31日，公司拥有境内授权发明专利186项，其中应用于主营业务并能够产业化的发明专利118项，满足条件 \\ \hline
最近三年营业收入复合增长率不低于25\%，或最近一年营业收入金额不低于3亿元 & \mbox{\kxv 是\hspace{0.4\ccwd}\kx 否} & 公司2023—2025年度营业收入复合增长率为26.81\%；2025年度营业收入为132,594.18万元，均满足条件 \\ \hline
\end{longtable}
\end{gstab}

% 未盈利等适用《科创属性评价指引》例外情形（标准二：国家级奖项/重大专项/
% 进口替代/50项以上发明专利等）的发行人，参照上表格式增列「科创属性评价
% 标准二」对照表。

\subsection{发行人报告期主要财务数据和财务指标}

公司报告期财务报表已经示例会计师审计，并由其出具了标准无保留意见的《审计报
告》（示例会所审字〔2026〕第××号）。报告期主要财务数据及财务指标如下：

\biaodw{万元}
\begin{gstab}
\begin{longtable}{|L{40mm}|R{32mm}|R{32mm}|R{32mm}|}
\hline
\thA{项目} & \thBw{32mm}{2025年度／\newline 2025年12月31日} & \thBw{32mm}{2024年度／\newline 2024年12月31日} & \thBw{32mm}{2023年度／\newline 2023年12月31日} \\ \hline
\endfirsthead
\hline
\thA{项目} & \thBw{32mm}{2025年度／\newline 2025年12月31日} & \thBw{32mm}{2024年度／\newline 2024年12月31日} & \thBw{32mm}{2023年度／\newline 2023年12月31日} \\ \hline
\endhead
资产总额 & 168,392.45 & 128,475.36 & 96,384.52 \\ \hline
归属于母公司股东的权益 & 97,055.27 & 75,078.92 & 58,236.41 \\ \hline
营业收入 & 132,594.18 & 105,238.47 & 82,456.32 \\ \hline
归属于母公司股东的净利润 & 21,976.35 & 16,842.51 & 12,035.28 \\ \hline
扣除非经常性损益后归属于母公司股东的净利润 & 21,235.68 & 16,238.47 & 11,562.83 \\ \hline
经营活动产生的现金流量净额 & 19,435.82 & 14,528.36 & 10,236.45 \\ \hline
综合毛利率 & 44.02\% & 43.18\% & 42.35\% \\ \hline
研发投入占营业收入的比例 & 9.72\% & 9.34\% & 8.99\% \\ \hline
基本每股收益（元／股） & 0.61 & 0.47 & 0.33 \\ \hline
加权平均净资产收益率 & 25.53\% & 25.27\% & 23.05\% \\ \hline
资产负债率（合并） & 42.36\% & 41.56\% & 39.58\% \\ \hline
\end{longtable}
\end{gstab}

\subsection{发行人选择的具体上市标准}

公司选择适用《上海证券交易所科创板股票上市规则》第2.1.2条第（一）项规定的
上市标准：``预计市值不低于人民币10亿元，最近两年净利润均为正且累计净利润不
低于人民币5,000万元''。公司2024年度、2025年度扣除非经常性损益前后孰低的净
利润分别为16,238.47万元和21,235.68万元，均为正且累计为37,474.15万元；结合
可比公司估值水平及前次外部股权融资估值，公司预计发行后总市值不低于人民币10
亿元，满足所选上市标准。

\subsection{公司治理特殊安排等重要事项}

截至本招股说明书签署日，公司不存在特别表决权股份、协议控制（VIE）架构或者
其他类似特殊安排。根据《公司法》及《公司章程》的规定，公司不设监事会，由董
事会审计委员会行使《公司法》规定的监事会职权。

\subsection{募集资金主要用途}

本次发行募集资金扣除发行费用后，将按轻重缓急顺序投资于以下项目：

\biaodw{万元}
\begin{gstab}
\begin{longtable}{|C{12mm}|L{58mm}|R{33mm}|R{33mm}|}
\hline
\thA{序号} & \thB{项目名称} & \thB{投资总额} & \thB{拟使用募集资金} \\ \hline
\endfirsthead
\hline
\thA{序号} & \thB{项目名称} & \thB{投资总额} & \thB{拟使用募集资金} \\ \hline
\endhead
1 & 新一代大模型研发及训练平台建设项目 & 82,456.00 & 82,456.00 \\ \hline
2 & 智能体应用与行业解决方案研发项目 & 35,208.00 & 35,208.00 \\ \hline
3 & 补充流动资金 & 32,336.00 & 32,336.00 \\ \hline
\multicolumn{2}{|c|}{\bfseries 合计} & {\bfseries 150,000.00} & {\bfseries 150,000.00} \\ \hline
\end{longtable}
\end{gstab}

若本次发行实际募集资金净额少于上述项目拟使用募集资金总额，不足部分由公司自
筹解决。募集资金具体运用情况详见本招股说明书``第七节\hspace{0.5\ccwd}募集资
金运用与未来发展规划''。
